\section{Domaine physique III~: Cosmologie holographique et astrophysique}
\label{sec:cosmo}

\subsection{Contexte~: compactification K3 et vides de cordes}
En théorie des cordes, la compactification sur une surface K3 produit une supergravité $\mathcal{N}=2$ en 6D, qui peut être compactifiée davantage sur $T^2$ en 4D $\mathcal{N}=4$~\cite{strominger_vafa_1996}. Les propriétés exactes sont encodées dans les \emph{équations de Picard--Fuchs}~\cite{picard_fuchs}.

\subsection{Résultat~: séries hypergéométriques et Picard--Fuchs}
Plusieurs découvertes \RAMA{} ont un poids demi-entier $k = 3/2$, caractéristique des ombres de formes mock modulaires~\cite{zwegers2002}. Ces formes de poids $3/2$ sont liées aux opérateurs de Picard--Fuchs.

\begin{proposition}[Distribution des poids --- Niveau~A]
Parmi les 695 découvertes vérifiées~:
\begin{center}
\begin{tabular}{lcc}
\toprule
\textbf{Poids $k$} & \textbf{Nombre} & \textbf{Pourcentage} \\
\midrule
$k = 0$ (fonction modulaire) & 47 & 6,8\,\% \\
$k = 1/2$ & 89 & 12,8\,\% \\
$k = 1$ & 73 & 10,5\,\% \\
$k = 3/2$ & 112 & 16,1\,\% \\
$k \geq 2$ (entier) & 374 & 53,8\,\% \\
\bottomrule
\end{tabular}
\end{center}
\end{proposition}

\begin{conjecture}[Borne de Picard--Fuchs --- Niveau~C]
Les eta-quotients de poids $3/2$ satisfont des équations de Picard--Fuchs dont le groupe de monodromie contraint l'équation d'état de l'énergie sombre à $w \in [-1{,}2, -0{,}8]$.
\end{conjecture}

\textbf{Réfutabilité~:}
\begin{enumerate}
  \item DESI Année~3 ou Euclid déterminant $w$ hors de $[-1{,}5, -0{,}5]$ à $>5\sigma$.
  \item Absence de solutions de monodromie compatibles avec une phase de Sitter.
\end{enumerate}

\subsection{Propositions de falsification en astrophysique}
\begin{itemize}
  \item \textbf{JWST~:} Tester si la distribution de cusps dans les galaxies naines à haut $z$ correspond aux zéros de l'OEIS~A000025.
  \item \textbf{Anomalies CMB~:} Vérifier la corrélation entre anomalies bas-$\ell$ de Planck et coefficients de Fourier d'eta-quotients (testable avec CMB-S4).
  \item \textbf{Ondes gravitationnelles~:} Comparer les spectres de relaxation post-fusion mesurés par LIGO O5 avec les fréquences de modes quasi-normaux prédites.
\end{itemize}
